\documentclass[10pt,a4paper]{amsart}
\usepackage{amsmath,amssymb,amsthm,graphicx}
\usepackage{amsmath}
\usepackage{amssymb}
\usepackage{hyperref}
\newtheorem{Theorem}{Theorem}[section]
\newtheorem{Corollary}[Theorem]{Corollary}
\newtheorem{Lemma}[Theorem]{Lemma}
\newtheorem{Proposition}[Theorem]{Proposition}
{\theoremstyle{definition}
\newtheorem{Note}[Theorem]{Note}
\newtheorem{Definition}[Theorem]{Definition}
\newtheorem{Problem}[Theorem]{Problem}
}
\usepackage{cleveref}
\crefname{originaltheorem}{original theorem}{original theorems}
\Crefname{originaltheorem}{Original theorem}{Original theorems}
\makeatletter
\newlabel{prop:altpres}{{1.6}{source}{}{sourcepointer-1}{}}
\newlabel{prop:altpres@cref}{{[originaltheorem][0][]1.6}{[1][0][]source}{}{}{}}
\makeatother
\title{Three inclusions describing the commutator of the algebra O in terms of the ideal generated by AB minus BA}
\author{Paul Terwilliger}
\date{doi:10.3842/SIGMA.2011.069; CC BY4.0}
\begin{document}
\maketitle
\noindent\textit{Editorial note.} All delivered mathematical text is the author's own, copied from the pinned publisher source sources/sigma-2011-069-askey-wilson/pinned-original-source.tex; no statement or proof step was rewritten, reordered or removed.\par\medskip
\noindent\textit{Editorial note.} Preamble provenance: the retained preamble is the source's own layout adapter plus the author's macro block at lines 10-18 of sources/sigma-2011-069-askey-wilson/pinned-original-source.tex (everything between the publisher class metadata and \textbackslash{}begin\{document\}, comments excluded). Nothing was added, no macro was invented, and the publisher document class is deliberately not reproduced; the excerpt serves only as preamble, never as a span coordinate.\par\medskip
\noindent\textit{Editorial note.} External references are navigation pointers only; every deliberately kept step of the author's argument stays inside the counted spans. The displayed numbers of those pointers are the author's own numbering under the author's \textbackslash{}newtheorem declarations (the source has no subsections, so numbering is section.index), computed mechanically by the lane's scanner.\par\medskip
\noindent\textit{Editorial note.} The single reference to another labelled item is Theorem prop:altpres, recorded as a navigation stub of kind theorem carrying the author's own number; the computation that consumes it stays inside the counted span.\par\medskip
\noindent\textit{Editorial note.} Clause (iii) is closed by the author in one line from clauses (i) and (ii) together with centrality of gamma; that sentence is kept verbatim.\par\medskip
\noindent\textit{Editorial note.} Licence binding: the card's source\_version now carries the value recorded in this source's own licence evidence (sources/sigma-2011-069-askey-wilson/license-observation.json), namely 'SIGMA7(2011),069; publisher final form2011-07-09'. The adopt step of the publication line compares the two and refused the earlier revision; the card is aligned to the evidence rather than the evidence to the card.\par\medskip
\noindent\textit{Allowed standard prerequisites (editorial).}\begin{itemize}
\item The algebra O is generated by A and B. Same source, set up before the statement; used inside the counted proof at lines 3070-3073. Makes the quotient O/R commutative, which is the whole content of clause (i).
\item The alternating presentation relation of the algebra, third displayed equation. Same source, Theorem labelled prop:altpres; invoked inside the counted proof at lines 3081-3085. Identifies the right-hand side of that relation with a nonzero scalar multiple of [A,B] gamma, which gives clause (ii).
\end{itemize}\medskip
\ifdefined\hypertarget\hypertarget{sourcepointer-1}{}\fi
\noindent\textit{Source reference.} Original theorem 1.6: sources/sigma-2011-069-askey-wilson/pinned-original-source.tex, source lines 3082--3082. Navigation pointer for \textbackslash{}ref\{prop:altpres\} used inside the counted span (line 3082); the argument that consumes it is retained in the span.\par\medskip
\begin{Lemma}
\label{lem:three}
The following $(i)$--$(iii)$ hold.
\begin{enumerate}\itemsep=0pt
\item[$(i)$]
$\lbrack \mathbb O,\mathbb O\rbrack \subseteq
\mathbb O \lbrack A,B\rbrack \mathbb O$.
\item[$(ii)$]
$  \lbrack A,B\rbrack
\gamma \in
\lbrack \mathbb O,\mathbb O\rbrack $.
\item[$(iii)$]
$\mathbb O \lbrack A,B\rbrack \mathbb O \gamma \subseteq
\mathbb O \lbrack A,B\rbrack \mathbb O$.
\end{enumerate}
\end{Lemma}

\begin{proof}\sloppy
$(i)$
Abbreviate $R=
\mathbb O \lbrack A,B\rbrack \mathbb O$ and consider the
quotient algebra $\mathbb O/R$.
The ele\-ments~$A$,~$B$ generate~$\mathbb O$, and
these generators satisfy
$\lbrack A,B\rbrack \in R$.
Therefore $A+R$, $B+R$ generate
$\mathbb O/R$, and these generators commute.
This shows that $\mathbb O/R$ is commutative.
Consider the canonical map $\mathbb O \to \mathbb O/R$.
This map has kernel~$R$. The map sends
$\lbrack \mathbb O,\mathbb O\rbrack \mapsto 0 $ since~$\mathbb O/R$ is commutative. Therefore
$\lbrack \mathbb O,\mathbb O\rbrack \subseteq R$.

  $(ii)$ In the third displayed equation of
Theorem~\ref{prop:altpres},
the expression on the right is a nonzero
scalar multiple of
$\lbrack A,B \rbrack \gamma $.
The expression on the left
is equal to
$\lbrack A^2,B^2 \rbrack +(q^2+q^{-2})
\lbrack B,ABA \rbrack $
and is therefore in
$\lbrack \mathbb O,\mathbb O\rbrack $.
The result follows.

  $(iii)$ By $(i)$, $(ii)$ above and
since $\gamma$ is central.
\end{proof}

\begin{thebibliography}{99}

\footnotesize\itemsep=0pt
\bibitem{alpern}
Alperin R.C.,
Notes: ${\rm  {PSL}}_2(\mathbb Z)= \mathbb Z_2 \star \mathbb Z_3$,
\href{http://dx.doi.org/10.2307/2324963}{{\it Amer. Math. Monthly}} {\bf 100} (1993), 385--386.
\end{thebibliography}
\end{document}
